\chapter{Complexity — Alphabet Reduction}
%%% generated by the blueprint pipeline from TCSlib.Complexity.TuringMachine.Robustness.AlphabetReduction
\section{Overview}

This module proves the in-model analogue of [AB09, Claim 1.5]: a Turing machine over any
finite alphabet $\Gamma$ (with a read-only input tape, a write-only output tape, and $k$
work tapes) that computes a binary string function via a bit embedding
$e : \{0,1\} \hookrightarrow \Gamma$ within time $T$ is simulated by a machine over the
binary alphabet with the \emph{same} number of work tapes within time $c \cdot (T(n)+1)$
for a constant $c$ depending only on $\abs{\Gamma}$. Each work-tape symbol is stored as a
fixed-width block of $N+1$ optional bits via an injective code (the implementation uses a
one-hot code of width $\abs{\Gamma}+1$, with logical blank coded by the all-blank block);
one source step is simulated by a rightward read sweep, a dispatch transition, a leftward
write sweep, and a whole-block head move, for exactly $3(N+1)+2$ simulator steps per
source step.

\section{Declarations}

\begin{definition}[Block of optional bits]
\lean{Turing.FinTM.ArBlock}
\label{Turing.FinTM.ArBlock}
\leanok
For $N \in \mathbb{N}$, a block is a vector of $N+1$ cells indexed by
$\{0,\dots,N\}$, each cell holding either a bit or the blank. Blocks serve as the
fixed-width binary encodings of single work-tape symbols.
\end{definition}

\begin{definition}[Block index of a tape position]
\lean{Turing.FinTM.arCell}
\label{Turing.FinTM.arCell}
\leanok
For $N \in \mathbb{N}$ and an integer tape position $p \in \mathbb{Z}$, the index
$\lfloor p/(N+1) \rfloor$ of the length-$(N+1)$ block of cells containing $p$.
\end{definition}

\begin{definition}[Offset of a tape position within its block]
\lean{Turing.FinTM.arIndex}
\label{Turing.FinTM.arIndex}
\leanok
For $N \in \mathbb{N}$ and $p \in \mathbb{Z}$, the nonnegative remainder
$p \bmod (N+1)$, viewed as an element of $\{0,\dots,N\}$: the offset of position $p$
inside the length-$(N+1)$ block containing it.
\end{definition}

\begin{definition}[Tape position of a block coordinate]
\lean{Turing.FinTM.arPos}
\label{Turing.FinTM.arPos}
\leanok
For $N \in \mathbb{N}$, a block index $z \in \mathbb{Z}$, and an offset
$j \in \mathbb{N}$, the integer tape position $(N+1)z + j$ of the $j$-th cell of the
$z$-th length-$(N+1)$ block.
\end{definition}

\begin{lemma}[Coordinates of a combined position]
\lean{Turing.FinTM.arCoords}
\label{Turing.FinTM.arCoords}
\leanok
\uses{Turing.FinTM.arCell, Turing.FinTM.arIndex, Turing.FinTM.arPos}
\difficulty{3}
For every $N \in \mathbb{N}$, every $z \in \mathbb{Z}$, and every
$j \in \{0,\dots,N\}$, the position $p = (N+1)z + j$ satisfies
$\lfloor p/(N+1) \rfloor = z$ and $p \bmod (N+1) = j$.
\end{lemma}

\begin{lemma}[Recombining the coordinates of a position]
\lean{Turing.FinTM.arPos_coords}
\label{Turing.FinTM.arPos_coords}
\leanok
\uses{Turing.FinTM.arCell, Turing.FinTM.arIndex, Turing.FinTM.arPos}
\difficulty{3}
For every $N \in \mathbb{N}$ and every $p \in \mathbb{Z}$,
\[
(N+1)\left\lfloor \frac{p}{N+1} \right\rfloor + \left(p \bmod (N+1)\right) = p .
\]
\end{lemma}

\begin{lemma}[Injectivity of block coordinates]
\lean{Turing.FinTM.arPos_eq}
\label{Turing.FinTM.arPos_eq}
\leanok
\uses{Turing.FinTM.arCell, Turing.FinTM.arCoords, Turing.FinTM.arIndex, Turing.FinTM.arPos}
\difficulty{3}
For $N \in \mathbb{N}$, $z, z' \in \mathbb{Z}$, and $j, j' \in \{0,\dots,N\}$, one has
$(N+1)z + j = (N+1)z' + j'$ if and only if $z = z'$ and $j = j'$.
\end{lemma}

\begin{definition}[Binary tape induced by a family of blocks]
\lean{Turing.FinTM.arTape}
\label{Turing.FinTM.arTape}
\leanok
\uses{Turing.FinTM.ArBlock, Turing.FinTM.arCell, Turing.FinTM.arIndex}
Given $N \in \mathbb{N}$ and a family $F$ assigning to every integer $z$ a block of
$N+1$ optional bits, the induced binary tape $\mathbb{Z} \to \{0,1\} \cup
\{\text{blank}\}$: position $p$ holds entry $p \bmod (N+1)$ of the block
$F(\lfloor p/(N+1) \rfloor)$.
\end{definition}

\begin{lemma}[Reading the induced tape at a block coordinate]
\lean{Turing.FinTM.arTape_at}
\label{Turing.FinTM.arTape_at}
\leanok
\uses{Turing.FinTM.arCoords, Turing.FinTM.arPos, Turing.FinTM.arTape}
\difficulty{1}
Let $F$ assign to every integer a block of $N+1$ optional bits. The binary tape induced
by $F$ (position $p$ holding entry $p \bmod (N+1)$ of block
$F(\lfloor p/(N+1)\rfloor)$) reads, at position $(N+1)z + j$ with $z \in \mathbb{Z}$
and $j \in \{0,\dots,N\}$, exactly the entry $j$ of $F(z)$.
\end{lemma}

\begin{lemma}[Extensionality along block coordinates]
\lean{Turing.FinTM.arTape_ext}
\label{Turing.FinTM.arTape_ext}
\leanok
\uses{Turing.FinTM.arCell, Turing.FinTM.arIndex, Turing.FinTM.arPos, Turing.FinTM.arPos_coords}
\difficulty{1}
Fix $N \in \mathbb{N}$. Two binary tapes $u, v : \mathbb{Z} \to \{0,1\} \cup
\{\text{blank}\}$ that agree at every position of the form $(N+1)z + j$ with
$z \in \mathbb{Z}$ and $j \in \{0,\dots,N\}$ are equal.
\end{lemma}

\begin{lemma}[Pointwise update of an induced tape]
\lean{Turing.FinTM.arTape_update}
\label{Turing.FinTM.arTape_update}
\leanok
\uses{Turing.FinTM.arPos, Turing.FinTM.arPos_eq, Turing.FinTM.arTape, Turing.FinTM.arTape_at, Turing.FinTM.arTape_ext}
\difficulty{2}
Let $F$ assign to every integer a block of $N+1$ optional bits, let $z \in \mathbb{Z}$,
$j \in \{0,\dots,N\}$, and let $b$ be an optional bit. Overwriting the binary tape
induced by $F$ at position $(N+1)z + j$ with the value $b$ yields the binary tape
induced by the family that agrees with $F$ everywhere except that entry $j$ of block
$z$ becomes $b$.
\end{lemma}

\begin{definition}[One-hot block code]
\lean{Turing.FinTM.arCode}
\label{Turing.FinTM.arCode}
\leanok
\uses{Turing.FinTM.ArBlock}
For a finite alphabet $\Gamma$, an injection of $\Gamma \cup \{\text{blank}\}$ into
blocks of $\abs{\Gamma}+1$ optional bits: the blank maps to the all-blank block, and a
symbol $x \in \Gamma$ maps to the block whose entry at offset $j$ is the genuine bit
that is $1$ exactly when $j$ is the index of $x$ under a fixed enumeration of
$\Gamma$.
\end{definition}

\begin{definition}[Block decoder]
\lean{Turing.FinTM.arDecode}
\label{Turing.FinTM.arDecode}
\leanok
\uses{Turing.FinTM.ArBlock}
Given an injection $E$ of $\Gamma \cup \{\text{blank}\}$ into blocks of $N+1$ optional
bits, a total decoding map from blocks to $\Gamma \cup \{\text{blank}\}$ chosen as a
left inverse of $E$: every block in the image of $E$ is sent to its unique preimage.
\end{definition}

\begin{lemma}[Decoding inverts coding]
\lean{Turing.FinTM.arDecode_code}
\label{Turing.FinTM.arDecode_code}
\leanok
\uses{Turing.FinTM.ArBlock, Turing.FinTM.arDecode}
\difficulty{1}
Let $E$ be an injection of $\Gamma \cup \{\text{blank}\}$ into blocks of $N+1$ optional
bits, and let the decoder be the associated left inverse of $E$. Then for every
$a \in \Gamma \cup \{\text{blank}\}$, decoding the block $E(a)$ returns $a$.
\end{lemma}

\begin{definition}[Total bit decoder]
\lean{Turing.FinTM.arBit}
\label{Turing.FinTM.arBit}
\leanok
For an embedding $e : \{0,1\} \hookrightarrow \Gamma$ of the two bits into an alphabet
$\Gamma$ and a symbol $a \in \Gamma$, the bit that is $1$ exactly when
$a = e(1)$; in particular every symbol outside the image of $e$ decodes to $0$.
\end{definition}

\begin{lemma}[Bit decoding inverts the bit embedding]
\lean{Turing.FinTM.arBit_embed}
\label{Turing.FinTM.arBit_embed}
\leanok
\uses{Turing.FinTM.arBit}
\difficulty{2}
Let $e : \{0,1\} \hookrightarrow \Gamma$ be an embedding of the bits into an alphabet
$\Gamma$. For every bit $b$, applying the total bit decoder (which sends $a \in \Gamma$
to $1$ exactly when $a = e(1)$) to $e(b)$ returns $b$.
\end{lemma}

\begin{definition}[Pending transition record]
\lean{Turing.FinTM.ArPending}
\label{Turing.FinTM.ArPending}
\leanok
For an alphabet $\Gamma$, a state type $S$, and $k \in \mathbb{N}$ work tapes, a
pending record consists of an optional next state (absence meaning halt), for each of
the $k$ work tapes an element of $\Gamma \cup \{\text{blank}\}$ to be written, and for
each work tape a head movement direction in $\{-1,0,+1\}$.
\end{definition}

\begin{definition}[Simulator state space]
\lean{Turing.FinTM.ArState}
\label{Turing.FinTM.ArState}
\leanok
\uses{Turing.FinTM.ArBlock, Turing.FinTM.ArPending}
The finite state space of the block-code simulator with parameters $\Gamma$, $S$, $k$,
$N$: either a reading state, consisting of a source state $q \in S$, a counter
$j \in \{0,\dots,N+1\}$, and for each of the $k$ work tapes a partially filled buffer
block of $N+1$ optional bits; or a pending record (optional next state, per-tape
symbol to write, per-tape direction) together with either a write counter in
$\{0,\dots,N\}$ or a move counter in $\{0,\dots,N+1\}$.
\end{definition}

\begin{definition}[Initial reading state]
\lean{Turing.FinTM.arReady}
\label{Turing.FinTM.arReady}
\leanok
\uses{Turing.FinTM.ArState}
For a source state $q$, the simulator state at a macro-step boundary: the reading state
with source state $q$, counter $0$, and all-blank buffer blocks.
\end{definition}

\begin{definition}[Pending record of a source action]
\lean{Turing.FinTM.arPending}
\label{Turing.FinTM.arPending}
\leanok
\uses{Turing.Action, Turing.FinTM.ArPending}
Given the vector of symbols currently scanned under the $k$ work heads and a source
machine action $a$ (next state, per-tape optional write and head move, input-head move,
optional emission), the pending record retaining: the next state of $a$; for each work
tape the effective new symbol (the symbol $a$ writes, or the scanned symbol when $a$
writes nothing); and the head movement direction of $a$ on each work tape.
\end{definition}

\begin{definition}[The binary block-code simulator machine]
\lean{Turing.FinTM.arTM}
\label{Turing.FinTM.arTM}
\leanok
\uses{Turing.FinTM, Turing.FinTM.ArBlock, Turing.FinTM.ArState, Turing.FinTM.arBit, Turing.FinTM.arDecode, Turing.FinTM.arPending, Turing.FinTM.arReady}
Given a finite alphabet $\Gamma$, an injective block code $E$ of
$\Gamma \cup \{\text{blank}\}$ into blocks of $N+1$ optional bits, a bit embedding
$e : \{0,1\} \hookrightarrow \Gamma$, and a Turing machine $M$ over $\Gamma$ (read-only
input tape, write-only output tape, $k$ work tapes), the simulating machine over the
binary alphabet with the same number $k$ of work tapes. Its control cycles through:
$N+1$ reading steps that copy the block under each work head into a buffer while moving
right; one dispatch step that decodes the buffers, applies the transition of $M$ (with
the scanned input bit mapped through $e$), performs the input-head move of $M$, emits
the bit-decoded image of any emitted symbol, stores the resulting pending record, and
moves each work head one cell left; $N+1$ writing steps that write the coded new
symbols right to left; and $N+2$ moving steps that shift each work head one whole block
in its source direction and then re-enter the reading state (halting if the source
transition halts).
\end{definition}

\begin{definition}[Encoded configuration]
\lean{Turing.FinTM.arCfg}
\label{Turing.FinTM.arCfg}
\leanok
\uses{Turing.Cfg, Turing.FinTM.ArBlock, Turing.FinTM.ArState, Turing.FinTM.arBit, Turing.FinTM.arPos, Turing.FinTM.arReady, Turing.FinTM.arTape}
Let $E$ be an injective block code of $\Gamma \cup \{\text{blank}\}$ into blocks of
$N+1$ optional bits, $e : \{0,1\} \hookrightarrow \Gamma$ a bit embedding, $x$ a binary
string, and $c$ a configuration of a $\Gamma$-machine with $k$ work tapes on the
letterwise $e$-image of $x$. The encoded simulator configuration on the genuine binary
input $x$: its state is the initial reading state for the state of $c$ (halted when $c$
is halted); the input-head position is unchanged; work tape $i$ is the binary tape
induced by coding the corresponding tape of $c$ blockwise through $E$; work head $i$
sits at offset $0$ of the block of its source head position; and the output of $c$ is
decoded bitwise (a symbol $a$ becoming $1$ exactly when $a = e(1)$).
\end{definition}

\begin{lemma}[Scanned input symbol of the encoded configuration]
\lean{Turing.FinTM.arCfg_input}
\label{Turing.FinTM.arCfg_input}
\leanok
\uses{Turing.Cfg, Turing.Cfg.inputSymbol, Turing.FinTM.ArBlock, Turing.FinTM.arCfg}
\difficulty{3}
Let $e : \{0,1\} \hookrightarrow \Gamma$ be a bit embedding, $x$ a binary string, and
$c$ a configuration of a $\Gamma$-machine on the letterwise $e$-image of $x$. The input
symbol scanned by the encoded simulator configuration of $c$ (which reads the genuine
binary input $x$ at the same input-head position), mapped through $e$, equals the input
symbol scanned by $c$; in particular both are absent for the same head positions.
\end{lemma}

\begin{definition}[Partially filled read buffer]
\lean{Turing.FinTM.arBuffer}
\label{Turing.FinTM.arBuffer}
\leanok
\uses{Turing.FinTM.ArBlock}
Given an injective block code $E$ of $\Gamma \cup \{\text{blank}\}$ into blocks of
$N+1$ optional bits, a vector assigning to each of $k$ work tapes a scanned element of
$\Gamma \cup \{\text{blank}\}$, and a count $j \in \mathbb{N}$: the per-tape buffer
after $j$ reading steps, namely for each tape $i$ the block whose entries at offsets
$< j$ agree with the code of the scanned symbol of tape $i$ and whose remaining entries
are blank.
\end{definition}

\begin{lemma}[Empty read buffer]
\lean{Turing.FinTM.arBuffer_zero}
\label{Turing.FinTM.arBuffer_zero}
\leanok
\uses{Turing.FinTM.ArBlock, Turing.FinTM.arBuffer}
\difficulty{1}
After zero reading steps the per-tape buffer (entries at offsets below the count taken
from the coded scanned symbols, blank elsewhere) is the all-blank block on every tape.
\end{lemma}

\begin{lemma}[Full read buffer]
\lean{Turing.FinTM.arBuffer_full}
\label{Turing.FinTM.arBuffer_full}
\leanok
\uses{Turing.FinTM.ArBlock, Turing.FinTM.arBuffer}
\difficulty{1}
After $N+1$ reading steps the per-tape buffer holds, on each tape $i$, the complete
code of the symbol scanned on tape $i$.
\end{lemma}

\begin{lemma}[One buffer fill step]
\lean{Turing.FinTM.arBuffer_update}
\label{Turing.FinTM.arBuffer_update}
\leanok
\uses{Turing.FinTM.ArBlock, Turing.FinTM.arBuffer}
\difficulty{4}
For $j \in \{0,\dots,N\}$, overwriting entry $j$ of the $j$-step buffer of each tape
with entry $j$ of the code of that tape's scanned symbol yields the $(j+1)$-step
buffer.
\end{lemma}

\begin{definition}[Configuration during the read sweep]
\lean{Turing.FinTM.arReadCfg}
\label{Turing.FinTM.arReadCfg}
\leanok
\uses{Turing.Cfg, Turing.Cfg.workTapeSymbols, Turing.FinTM.ArBlock, Turing.FinTM.ArState, Turing.FinTM.arBuffer, Turing.FinTM.arCfg, Turing.FinTM.arPos}
The simulator configuration $j$ steps into the read sweep of a macro step, for a source
configuration $c$, a source state $q$, and $j \in \{0,\dots,N+1\}$: it agrees with the
encoded configuration of $c$ except that its state is the reading state with source
state $q$, counter $j$, and the $j$-step buffers of the symbols scanned by $c$, and
each work head sits at offset $j$ of the block of its source head position.
\end{definition}

\begin{lemma}[Read sweep at counter zero]
\lean{Turing.FinTM.arReadCfg_zero}
\label{Turing.FinTM.arReadCfg_zero}
\leanok
\uses{Turing.Cfg, Turing.FinTM.ArBlock, Turing.FinTM.arBuffer_zero, Turing.FinTM.arCfg, Turing.FinTM.arReadCfg, Turing.FinTM.arReady}
\difficulty{1}
If the source configuration $c$ is in state $q$ (in particular not halted), then the
read-sweep configuration of $c$ at counter $0$ coincides with the encoded configuration
of $c$.
\end{lemma}

\begin{lemma}[Symbols scanned during the read sweep]
\lean{Turing.FinTM.arReadCfg_symbols}
\label{Turing.FinTM.arReadCfg_symbols}
\leanok
\uses{Turing.Cfg, Turing.Cfg.workTapeSymbols, Turing.FinTM.ArBlock, Turing.FinTM.arReadCfg, Turing.FinTM.arTape_at}
\difficulty{1}
For a counter $j \le N$, the work-tape symbols scanned by the read-sweep configuration
of a source configuration $c$ at counter $j$ are, on each tape $i$, the entry at offset
$j$ of the code of the symbol scanned by $c$ on tape $i$.
\end{lemma}

\begin{lemma}[One step of the read sweep]
\lean{Turing.FinTM.arReadCfg_step}
\label{Turing.FinTM.arReadCfg_step}
\leanok
\uses{Turing.Action, Turing.Action.apply, Turing.Cfg, Turing.Cfg.workTapeSymbols, Turing.FinTM, Turing.FinTM.ArBlock, Turing.FinTM.ArState, Turing.FinTM.arBuffer_update, Turing.FinTM.arCfg, Turing.FinTM.arPos, Turing.FinTM.arReadCfg, Turing.FinTM.arReadCfg_symbols, Turing.FinTM.arTM, Turing.MultiTapeTM.step}
\difficulty{4}
For a counter $j \le N$, one transition of the binary block-code simulator applied to
the read-sweep configuration of a source configuration $c$ with source state $q$ at
counter $j$ yields the read-sweep configuration at counter $j+1$: the scanned code bits
enter the buffers and every work head moves one cell right.
\end{lemma}

\begin{lemma}[Running the read sweep]
\lean{Turing.FinTM.arReadCfg_run}
\label{Turing.FinTM.arReadCfg_run}
\leanok
\uses{Turing.Cfg, Turing.FinTM, Turing.FinTM.ArBlock, Turing.FinTM.arCfg, Turing.FinTM.arReadCfg, Turing.FinTM.arReadCfg_step, Turing.FinTM.arReadCfg_zero, Turing.FinTM.arTM, Turing.MultiTapeTM.runFrom, Turing.MultiTapeTM.runFrom_succ_eq_step'}
\difficulty{4}
Let $c$ be a source configuration in state $q$ (not halted). For every $j \le N+1$,
running the binary block-code simulator for $j$ steps from the encoded configuration of
$c$ reaches the read-sweep configuration of $c$ at counter $j$.
\end{lemma}

\begin{definition}[Block family during the write sweep]
\lean{Turing.FinTM.arTail}
\label{Turing.FinTM.arTail}
\leanok
\uses{Turing.FinTM.ArBlock}
Given an injective block code $E$ of $\Gamma \cup \{\text{blank}\}$ into blocks of
$N+1$ optional bits, a source tape $t : \mathbb{Z} \to \Gamma \cup \{\text{blank}\}$, a
position $p \in \mathbb{Z}$, a symbol $w$ to be written at $p$, and a counter
$r \in \mathbb{N}$: the family of blocks coding the partially rewritten tape, where
every block $z \neq p$ is the code of $t(z)$, and block $p$ takes its entries at
offsets $\ge r$ from the code of $w$ and its entries at offsets $< r$ from the code of
$t(p)$.
\end{definition}

\begin{lemma}[Write-sweep family at full counter]
\lean{Turing.FinTM.arTail_full}
\label{Turing.FinTM.arTail_full}
\leanok
\uses{Turing.FinTM.ArBlock, Turing.FinTM.arTail}
\difficulty{2}
At counter $N+1$ (before anything is written) the partially rewritten block family for
tape $t$, write position $p$, and new symbol $w$ is simply the blockwise coding
$z \mapsto E(t(z))$ of the unmodified tape.
\end{lemma}

\begin{lemma}[Write-sweep family at counter zero]
\lean{Turing.FinTM.arTail_zero}
\label{Turing.FinTM.arTail_zero}
\leanok
\uses{Turing.FinTM.ArBlock, Turing.FinTM.arTail}
\difficulty{2}
At counter $0$ (after the write is complete) the partially rewritten block family for
tape $t$, write position $p$, and new symbol $w$ is the blockwise coding of the updated
tape that agrees with $t$ everywhere except that position $p$ now holds $w$.
\end{lemma}

\begin{lemma}[One leftward write extends the completed suffix]
\lean{Turing.FinTM.arTail_update}
\label{Turing.FinTM.arTail_update}
\leanok
\uses{Turing.FinTM.ArBlock, Turing.FinTM.arPos, Turing.FinTM.arPos_eq, Turing.FinTM.arTail, Turing.FinTM.arTape, Turing.FinTM.arTape_at, Turing.FinTM.arTape_ext}
\difficulty{4}
For $j \in \{0,\dots,N\}$, overwriting the binary tape induced by the counter-$(j+1)$
partially rewritten block family (tape $t$, write position $p$, new symbol $w$) at
tape position $(N+1)p + j$ with entry $j$ of the code of $w$ yields the binary tape
induced by the counter-$j$ family: blocks other than $p$ are unchanged, and within
block $p$ the only new entry is the one at offset $j$.
\end{lemma}

\begin{definition}[Configuration during the move sweep]
\lean{Turing.FinTM.arMoveCfg}
\label{Turing.FinTM.arMoveCfg}
\leanok
\uses{Turing.Action, Turing.Action.apply, Turing.Cfg, Turing.Cfg.workTapeSymbols, Turing.FinTM.ArBlock, Turing.FinTM.ArState, Turing.FinTM.arCfg, Turing.FinTM.arPending, Turing.FinTM.arPos}
The simulator configuration in the move sweep at counter $r \in \{0,\dots,N+1\}$, for a
source configuration $c$ and a source action $a$: the work tapes, input position, and
output are those of the encoded configuration of the result of applying $a$ to $c$; the
state carries the pending record of $a$ with move counter $r$; and work head $i$ sits
at position $(N+1)p_i + (N+1-r)\,d_i$, where $p_i$ is the source head position of $c$
on tape $i$ and $d_i \in \{-1,0,+1\}$ is the movement of $a$ on tape $i$.
\end{definition}

\begin{definition}[Configuration during the write sweep]
\lean{Turing.FinTM.arWriteCfg}
\label{Turing.FinTM.arWriteCfg}
\leanok
\uses{Turing.Action, Turing.Action.apply, Turing.Cfg, Turing.Cfg.workTapeSymbols, Turing.FinTM.ArBlock, Turing.FinTM.ArState, Turing.FinTM.arCfg, Turing.FinTM.arPending, Turing.FinTM.arPos, Turing.FinTM.arTail, Turing.FinTM.arTape}
The simulator configuration in the write sweep at counter $j \in \{0,\dots,N\}$, for a
source configuration $c$ and a source action $a$: the state carries the pending record
of $a$ with write counter $j$; work tape $i$ is the binary tape induced by the
partially rewritten block family in which, at the block of the source head position,
the entries at offsets $> j$ already hold the code of the effective new symbol (the
symbol $a$ writes, or the currently scanned symbol when $a$ writes nothing); work head
$i$ sits at offset $j$ of that block; and the input position and output are those of
the encoded configuration of the result of applying $a$ to $c$.
\end{definition}

\begin{lemma}[One step of the move sweep]
\lean{Turing.FinTM.arMoveCfg_step}
\label{Turing.FinTM.arMoveCfg_step}
\leanok
\uses{Turing.Action, Turing.Action.apply, Turing.Cfg, Turing.FinTM, Turing.FinTM.ArBlock, Turing.FinTM.ArState, Turing.FinTM.arCfg, Turing.FinTM.arMoveCfg, Turing.FinTM.arPending, Turing.FinTM.arTM, Turing.MultiTapeTM.step}
\difficulty{4}
For a counter $r \neq 0$, one transition of the binary block-code simulator applied to
the move-sweep configuration (source configuration $c$, source action $a$) at counter
$r$ yields the move-sweep configuration at counter $r-1$: each work head advances one
cell in its pending direction.
\end{lemma}

\begin{lemma}[Finishing the move sweep]
\lean{Turing.FinTM.arMoveCfg_finish}
\label{Turing.FinTM.arMoveCfg_finish}
\leanok
\uses{Turing.Action, Turing.Action.apply, Turing.Cfg, Turing.FinTM, Turing.FinTM.ArBlock, Turing.FinTM.arCfg, Turing.FinTM.arMoveCfg, Turing.FinTM.arPending, Turing.FinTM.arPos, Turing.FinTM.arTM, Turing.MultiTapeTM.step}
\difficulty{4}
One transition of the binary block-code simulator applied to the move-sweep
configuration (source configuration $c$, source action $a$) at counter $0$ yields the
encoded configuration of the result of applying $a$ to $c$; in particular the
simulator halts exactly when the next state of $a$ is absent.
\end{lemma}

\begin{lemma}[Running the move sweep]
\lean{Turing.FinTM.arMoveCfg_run}
\label{Turing.FinTM.arMoveCfg_run}
\leanok
\uses{Turing.Action, Turing.Action.apply, Turing.Cfg, Turing.FinTM, Turing.FinTM.ArBlock, Turing.FinTM.arCfg, Turing.FinTM.arMoveCfg, Turing.FinTM.arMoveCfg_finish, Turing.FinTM.arMoveCfg_step, Turing.FinTM.arTM, Turing.MultiTapeTM.runFrom, Turing.MultiTapeTM.runFrom_succ_eq_step}
\difficulty{4}
For every counter $r \le N+1$, running the binary block-code simulator for $r+1$ steps
from the move-sweep configuration (source configuration $c$, source action $a$) at
counter $r$ reaches the encoded configuration of the result of applying $a$ to $c$.
\end{lemma}

\begin{lemma}[One step of the write sweep]
\lean{Turing.FinTM.arWriteCfg_step}
\label{Turing.FinTM.arWriteCfg_step}
\leanok
\uses{Turing.Action, Turing.Action.apply, Turing.Cfg, Turing.FinTM, Turing.FinTM.ArBlock, Turing.FinTM.ArState, Turing.FinTM.arCfg, Turing.FinTM.arPending, Turing.FinTM.arPos, Turing.FinTM.arTM, Turing.FinTM.arTail_update, Turing.FinTM.arWriteCfg, Turing.MultiTapeTM.step}
\difficulty{4}
For a counter $j \neq 0$, one transition of the binary block-code simulator applied to
the write-sweep configuration (source configuration $c$, source action $a$) at counter
$j$ yields the write-sweep configuration at counter $j-1$: on each tape the entry at
offset $j$ of the code of the effective new symbol is written and the head moves one
cell left.
\end{lemma}

\begin{lemma}[Finishing the write sweep]
\lean{Turing.FinTM.arWriteCfg_finish}
\label{Turing.FinTM.arWriteCfg_finish}
\leanok
\uses{Turing.Action, Turing.Action.apply, Turing.Action.apply_workTapes, Turing.Cfg, Turing.Cfg.workTapeSymbols, Turing.FinTM, Turing.FinTM.ArBlock, Turing.FinTM.arCfg, Turing.FinTM.arMoveCfg, Turing.FinTM.arPending, Turing.FinTM.arPos, Turing.FinTM.arTM, Turing.FinTM.arTail_update, Turing.FinTM.arTail_zero, Turing.FinTM.arTape, Turing.FinTM.arWriteCfg, Turing.MultiTapeTM.step}
\difficulty{5}
One transition of the binary block-code simulator applied to the write-sweep
configuration (source configuration $c$, source action $a$) at counter $0$ writes the
last code entry in place and yields the move-sweep configuration at counter $N+1$: the
work tapes now hold the complete blockwise coding of the tapes of the result of
applying $a$ to $c$.
\end{lemma}

\begin{lemma}[Running the write sweep]
\lean{Turing.FinTM.arWriteCfg_run}
\label{Turing.FinTM.arWriteCfg_run}
\leanok
\uses{Turing.Action, Turing.Cfg, Turing.FinTM, Turing.FinTM.ArBlock, Turing.FinTM.arMoveCfg, Turing.FinTM.arTM, Turing.FinTM.arWriteCfg, Turing.FinTM.arWriteCfg_finish, Turing.FinTM.arWriteCfg_step, Turing.MultiTapeTM.runFrom, Turing.MultiTapeTM.runFrom_succ_eq_step}
\difficulty{4}
For every counter $j \le N$, running the binary block-code simulator for $j+1$ steps
from the write-sweep configuration (source configuration $c$, source action $a$) at
counter $j$ reaches the move-sweep configuration at counter $N+1$.
\end{lemma}

\begin{lemma}[Dispatch step between the read and write sweeps]
\lean{Turing.FinTM.arReadCfg_dispatch}
\label{Turing.FinTM.arReadCfg_dispatch}
\leanok
\uses{Turing.Action.apply, Turing.Cfg, Turing.Cfg.inputSymbol, Turing.Cfg.workTapeSymbols, Turing.FinTM, Turing.FinTM.ArBlock, Turing.FinTM.arBuffer_full, Turing.FinTM.arCfg, Turing.FinTM.arCfg_input, Turing.FinTM.arDecode_code, Turing.FinTM.arPos, Turing.FinTM.arReadCfg, Turing.FinTM.arTM, Turing.FinTM.arTail_full, Turing.FinTM.arWriteCfg, Turing.MultiTapeTM.step, Turing.moveInputPos}
\difficulty{5}
Let $M$ be a Turing machine over a finite alphabet $\Gamma$, $c$ a configuration of $M$
on the letterwise $e$-image of a binary string, and $q$ a source state. One transition
of the binary block-code simulator applied to the read-sweep configuration of $c$ at
the full counter $N+1$ (all blocks buffered) yields the write-sweep configuration at
counter $N$ for the action produced by the transition function of $M$ on the state $q$,
the input symbol scanned by $c$, and the work-tape symbols scanned by $c$: the buffers
are decoded, the input head moves and any emission is bit-decoded and emitted, and each
work head steps left onto the last entry of its block.
\end{lemma}

\begin{lemma}[One macro cycle simulates one source step]
\lean{Turing.FinTM.arCfg_cycle}
\label{Turing.FinTM.arCfg_cycle}
\leanok
\uses{Turing.Cfg, Turing.FinTM, Turing.FinTM.ArBlock, Turing.FinTM.arCfg, Turing.FinTM.arMoveCfg_run, Turing.FinTM.arReadCfg, Turing.FinTM.arReadCfg_dispatch, Turing.FinTM.arReadCfg_run, Turing.FinTM.arTM, Turing.FinTM.arWriteCfg_run, Turing.MultiTapeTM.runFrom, Turing.MultiTapeTM.runFrom_add, Turing.MultiTapeTM.runFrom_of_halt, Turing.MultiTapeTM.step, Turing.MultiTapeTM.step_of_halt}
\difficulty{5}
Let $M$ be a Turing machine over a finite alphabet $\Gamma$ and $c$ a configuration of
$M$ on the letterwise $e$-image of a binary string. Running the binary block-code
simulator for $3(N+1)+2$ steps from the encoded configuration of $c$ yields the
encoded configuration of the result of one step of $M$ from $c$; when $c$ is halted
both encoded configurations are fixed points, so the identity holds in that case as
well.
\end{lemma}

\begin{lemma}[Running the simulator for many macro cycles]
\lean{Turing.FinTM.arCfg_run}
\label{Turing.FinTM.arCfg_run}
\leanok
\uses{Turing.Cfg, Turing.FinTM, Turing.FinTM.ArBlock, Turing.FinTM.arCfg, Turing.FinTM.arCfg_cycle, Turing.FinTM.arTM, Turing.MultiTapeTM.runFrom, Turing.MultiTapeTM.runFrom_add, Turing.MultiTapeTM.runFrom_succ_eq_step', Turing.MultiTapeTM.runFrom_zero}
\difficulty{4}
Let $M$ be a Turing machine over a finite alphabet $\Gamma$ and $c$ a configuration of
$M$ on the letterwise $e$-image of a binary string. For every $t \in \mathbb{N}$,
running the binary block-code simulator for $(3(N+1)+2)\cdot t$ steps from the encoded
configuration of $c$ yields the encoded configuration of the configuration reached by
$M$ after $t$ steps from $c$.
\end{lemma}

\begin{lemma}[Encoding of the initial configuration]
\lean{Turing.FinTM.arCfg_init}
\label{Turing.FinTM.arCfg_init}
\leanok
\uses{Turing.Cfg.init, Turing.FinTM, Turing.FinTM.ArBlock, Turing.FinTM.arCfg, Turing.FinTM.arPos, Turing.FinTM.arTM, Turing.FinTM.arTape, Turing.MultiTapeTM.initCfg}
\difficulty{3}
Suppose the injective block code $E$ sends the blank to the all-blank block. Then for
every Turing machine $M$ over a finite alphabet $\Gamma$, every bit embedding
$e : \{0,1\} \hookrightarrow \Gamma$, and every binary string $x$, the initial
configuration of the binary block-code simulator on input $x$ equals the encoded
initial configuration of $M$ on the letterwise $e$-image of $x$.
\end{lemma}

\begin{lemma}[Emitted symbols lie in the bit embedding's image]
\lean{Turing.FinTM.arEmission_in_image}
\label{Turing.FinTM.arEmission_in_image}
\leanok
\uses{Turing.FinTM, Turing.FinTM.ComputesFunInTimeVia, Turing.FinTM.computesInTime_iff, Turing.MultiTapeTM.initCfg, Turing.MultiTapeTM.outputSymbol, Turing.MultiTapeTM.outputSymbol_of_halt, Turing.MultiTapeTM.output_prefix, Turing.MultiTapeTM.runFrom, Turing.MultiTapeTM.runFrom_add, Turing.MultiTapeTM.runFrom_of_halt, Turing.MultiTapeTM.runFrom_succ_eq_step', Turing.MultiTapeTM.step_output}
\difficulty{5}
Let $M$ be a Turing machine over an alphabet $\Gamma$ and $e : \{0,1\}
\hookrightarrow \Gamma$ a bit embedding, and suppose $M$ computes the string function
$f : \{0,1\}^* \to \{0,1\}^*$ via $e$ within time $T$: for every binary string $x$,
the run of $M$ on the letterwise $e$-image of $x$ halts within $T(\abs{x})$ steps with
output the letterwise $e$-image of $f(x)$. If, in the run of $M$ on the $e$-image of
some binary input $x$, the transition taken from the configuration reached after $t$
steps emits a symbol $a \in \Gamma$, then $a = e(b)$ for some bit $b$.
\end{lemma}

\begin{lemma}[The simulator computes with constant-factor slowdown]
\lean{Turing.FinTM.arTM_computes}
\label{Turing.FinTM.arTM_computes}
\leanok
\uses{Turing.FinTM, Turing.FinTM.ArBlock, Turing.FinTM.ComputesInTime, Turing.FinTM.arBit_embed, Turing.FinTM.arCfg, Turing.FinTM.arCfg_init, Turing.FinTM.arCfg_run, Turing.FinTM.arTM, Turing.FinTM.computesInTime_iff}
\difficulty{3}
Let $\Gamma$ be a finite alphabet, $E$ an injective block code of $\Gamma \cup
\{\text{blank}\}$ into blocks of $N+1$ optional bits that sends the blank to the
all-blank block, $e : \{0,1\} \hookrightarrow \Gamma$ a bit embedding, and $M$ a
Turing machine over $\Gamma$. If, started on the letterwise $e$-image of a binary
string $x$, $M$ halts within $t$ steps with output the letterwise $e$-image of a
binary string $w$, then the binary block-code simulator, started on the input $x$
itself, halts within $(3(N+1)+2)\cdot t$ steps with output $w$.
\end{lemma}

\begin{theorem}[Alphabet reduction]
\lean{Turing.FinTM.alphabet_reduction}
\label{Turing.FinTM.alphabet_reduction}
\leanok
\uses{Turing.FinTM, Turing.FinTM.ComputesFunInTime, Turing.FinTM.ComputesFunInTimeVia, Turing.FinTM.ComputesInTime.mono, Turing.FinTM.arCode, Turing.FinTM.arTM, Turing.FinTM.arTM_computes}
\difficulty{2}
[AB09, Claim 1.5]. Let $\Gamma$ be a finite alphabet, $e : \{0,1\} \hookrightarrow
\Gamma$ an embedding of the two bits, $M$ a Turing machine over $\Gamma$ (with a
read-only input tape, a write-only output tape, and $k$ work tapes), $f$ a function on
binary strings, and $T : \mathbb{N} \to \mathbb{N}$. Suppose $M$ computes $f$ via $e$
within time $T$: for every binary string $x$, the run of $M$ on the letterwise
$e$-image of $x$ halts within $T(\abs{x})$ steps with output the letterwise $e$-image
of $f(x)$. Then there exist a constant $c \in \mathbb{N}$ and a Turing machine $M'$
over the binary alphabet with the same number $k$ of work tapes such that $M'$
computes $f$ within time $n \mapsto c\,(T(n)+1)$.
\end{theorem}
